%%
%% This is file `chemobabel-ja.tex',
%% generated with the docstrip utility.
%%
%% The original source files were:
%%
%% chemobabel.dtx  (with options: `ja')
%% 
%% IMPORTANT NOTICE:
%% 
%% For the copyright see the source file.
%% 
%% Any modified versions of this file must be renamed
%% with new filenames distinct from chemobabel-ja.tex.
%% 
%% For distribution of the original source see the terms
%% for copying and modification in the file chemobabel.dtx.
%% 
%% This generated file may be distributed as long as the
%% original source files, as listed above, are part of the
%% same distribution. (The sources need not necessarily be
%% in the same archive or directory.)
%%
%% This document describes the usage of `chemobabel.sty',
%% a LaTeX package for generating chemical structural formulas
%% using Open Babel and Inkscape.
%%
%% Copyright 2014-2022 Acetaminophen (Hironobu YAMASHITA)
%%   Email   :  h.y.acetaminophen[a t]gmail.com
%%   GitHub  :  https://github.com/aminophen
%%   Blog    :  http://acetaminophen.hatenablog.com/
%%   Twitter :  @aminophen
%%
\documentclass[12pt]{ltjsarticle}
\usepackage[margin=20mm]{geometry}
\usepackage[english]{babel}
\usepackage[nofontspec,defaultsups]{newpxtext}
\usepackage[zerostyle=c,straightquotes]{newtxtt}
\usepackage{newpxmath}
\usepackage{ascmac}
\usepackage[librsvg]{chemobabel}
\usepackage[bookmarksnumbered=true,hidelinks,%
 pdftitle={User Manual for chemobabel.sty},%
 pdfauthor={Hironobu Yamashita (Acetaminophen)},%
 pdfsubject={TeX \& LaTeX Advent Calendar 2014},%
 pdfkeywords={chemistry,open babel,chemdraw,smiles}]{hyperref}
\usepackage{footnotebackref}
\makeatletter
\long\def\@makecaption#1#2{%
  \vskip\abovecaptionskip
  \sbox\@tempboxa{#1: #2}%
  \global\@minipagefalse
  \hbox to\hsize{\hfil\box\@tempboxa\hfil}%
  \vskip\belowcaptionskip}
\makeatother
\newcount\TestCount
\def\XyM{\ifnum\fam=-1\relax\fam=0\relax\fi\TestCount=\fam
  X\kern-.30em\smash{\raise.50ex\hbox{$\fam\TestCount\Upsilon$}}%
  \kern-.30em{M}}
\def\XyMTeX{\XyM\kern-.1em\TeX}
\def\Lpack#1{\textsf{#1}}
\def\Lcode#1{\texttt{#1}}
\def\Lmeta#1{\textit{#1}}
\def\Lver#1{\Lpack{chemobabel} #1}
\def\TikZ{Ti\textit{k}Z}
\title{\Lpack{chemobabel} \\[1ex]
  \normalsize --- Chemical Structures from MDL Molfiles,
  ChemDraw Files or SMILES Notations --- \\
  化学構造式をMOLファイルやChemDrawファイル，SMILES表記法から自動生成
  }
\author{Hironobu Yamashita (@aminophen)}
\begin{document}
\renewcommand{\figurename}{Fig.\,}
\pagenumbering{roman}
\maketitle

この文書では，\textbf{化学構造式を生成}して\LaTeX の文書中に挿入するための
パッケージである\Lpack{chemobabel.sty}の使い方を説明します。
構造式はOpen Babelの機能を利用することにより，MDL Molfile, ChemDrawファイル，
さらにはSMILES表記法を含む\textbf{さまざまな化学データ形式から生成}できます。
以下の構造式はこの方法によって出力されたものです\footnote{特に断りがない
限り，この文書で使用した版は\Lcode{\input{obversion.txt}}です。}。

For English documentation, please refer to chemobabel-en.pdf.

\begin{figure}[ht]
  \centering
  \chemobabel[width=65mm]{draw/Firefly luciferin.mol}{} \quad
  \chemobabel*[width=90mm]{draw/Brevetoxin A.mol}
  \caption{Firefly luciferin \& Brevetoxin A (from MDL Molfiles)}
\end{figure}

\begin{figure}[ht]
  \centering
  \chemobabel[width=80mm]{draw/ATP.cdx}{} \hspace{1cm}
  \chemobabel*[width=30mm]{draw/Glucose.cdx}
  \caption{ATP (Adenosine triphosphate) \& Glucose (from ChemDraw files)}
\end{figure}
\begin{center}\small
Note: with Open Babel 3.1.1 the chemical structure of ATP from the .cdx file
becomes wrong.\\ See \url{https://github.com/openbabel/openbabel/issues/338}
\end{center}

\begin{figure}[ht]
  \centering
  \smilesobabel[width=36mm]{CCO}{} \hspace{1cm}
  \smilesobabel*[width=40mm]{CC(C(=O)O)N} \hspace{1cm}
  \smilesobabel*[width=30mm]{C1[C@H](C)C[C@@H](O)C1}
  \caption{Ethanol, Alanine \& (1\textit{S},3\textit{S})-3-Methylcyclopentanol (from SMILES notations)}
\end{figure}

\begin{figure}[ht]
  \centering
  \smilesobabel*[width=30mm]{CC(=O)Nc1ccc(cc1)O} \hspace{1cm}
  \smilesobabel*[width=40mm]{C([C@@H]1[C@H]([C@@H]([C@H]([C@H](O1)O)O)O)O)O}
  \caption{Acetaminophen (Paracetamol) \& $\alpha$-\textsc{d}-Glucose (from SMILES notations)}
\end{figure}

\begin{figure}[ht]
\centering
このパッケージはGitHubで管理しています。\par
\url{https://github.com/aminophen/chemobabel}\par
問い合わせや改善案などはGitHubのIssuesへお願いします。
\end{figure}

\begin{figure}[ht]
\centering
CTANからも入手可能です →
\url{https://www.ctan.org/pkg/chemobabel}\par
\TeX\ LiveやMiK\TeX でも配布されています。
以下はCTANの短い紹介文です。\par
\begin{screen}
This package provides a way to convert and include
chemical structure graphics from various chemical formats,
such as ChemDraw files, MDL molfile or SMILES notations
using Open Babel.
To use this \LaTeX\ package, it is necessary to enable execution
of the following external commands via \Lcode{latex -shell-escape}.
\begin{itemize}
 \item \Lcode{obabel} (Open Babel)
 \item \Lcode{inkscape} or \Lcode{rsvg-convert}
   (for SVG -> PDF/EPS conversion)
 \item \Lcode{pdfcrop} or \Lcode{ps2eps}
   (optional; for cropping large margins of PDF/EPS)
\end{itemize}
\end{screen}
\end{figure}

\clearpage
\setcounter{tocdepth}{3}
\tableofcontents

\clearpage
\pagenumbering{arabic}

\section{はじめに}

\subsection{開発の動機}

ご存じのとおり，\LaTeX は世界中で使われています。
しかし，化学構造式を描くという点で考えると，
\LaTeX 文書中に構造式を挿入する手段は非常に限られています。
この目的を達成する確実かつ容易な方法が存在しないことは，化学分野における
多くの研究者が\LaTeX システムを使わない要因になっていると考えられます。

すでに\href{http://www.ctan.org/}{CTAN}からいくつかのパッケージが利用可能です：
\begin{itemize}
  \item \href{http://www.ctan.org/pkg/xymtex}{\XyMTeX}:
    化学構造式を描くためのパッケージ集
  \item \href{http://www.ctan.org/pkg/chemfig}{\Lpack{chemfig}}:
    \TikZ を利用して構造式を描画するパッケージ
\end{itemize}
これらのパッケージは確実で，あらゆる構造式を一貫性をもって描画することができます。
しかし，十分にその機能を活用できるようになるには相当の熟練を要します。

化学者向けの専用ツール
（\href{http://www.cambridgesoft.com/Ensemble_for_Chemistry/ChemDraw/}{ChemDraw}など）
でPDFやEPS形式で出力し，\verb|\includegraphics|によって
構造式を取り込むという手段もあります。
この場合，化学用のフォーマット（\verb|.cdx| や \verb|.mol| など）と
画像フォーマット（\verb|.pdf| や \verb|.eps| など）の両方で保存する必要があります。

新しいパッケージ \Lpack{chemobabel.sty} は，\LaTeX 文書中に化学構造式を
挿入する必要がある化学者に新たな選択肢を提供することでしょう。
この方法では，\href{http://openbabel.org/}{Open Babel}によって
SVG形式の構造式画像を生成し，
さらに\href{https://inkscape.org/ja/}{Inkscape}によってSVGから
PDFまたはEPS形式の画像に変換します。
どちらのプログラムもオープンソースかつクロスプラットフォームであり，
盛んに開発が行われています。

\paragraph{Open Babelについて}

Open Babelはさまざまな化学データを扱うために特別に設計されたツールです。
オープンな共同プロジェクトで，分子モデリング・化学・固体物性・生化学や
関連分野のデータを誰もが検索・変換・分析・保存できることを目指しています。
詳細は\href{http://openbabel.org/}{公式サイト}を参照してください。

\clearpage

\subsection{着想とアプローチ}


最初にOpen Babelを用いてSMILES表記から化学構造式を生成するという
Noel O'Blog の記事\cite{NOB1}, \cite{NOB2}とコメント\cite{JLA}を見つけました。
これらの投稿に示された手法は非常におもしろいと思いましたが，
すべてのプログラムが \verb|-shell-escape| オプションによって実行されて
しまうことがあまり好きではありませんでした。
それと同じころ，あるソース中の \verb|\includegraphics| を
別のソースに書き出すというマクロ\cite{OKU}を偶然みつけ，
このとき私は \verb|-shell-escape| が必要な最小のコードだけを抽出する
という着想を得ました。

この方法はChemDrawファイルを画像に変換する目的にも応用できると考えたのですが，
この発想を実現したプロジェクトが見当たらず，ただ一つ見つけたSourceforgeの
\href{http://chemdrawinlatex.sourceforge.net/}{Chemdraw in \LaTeX}は
すでに終了していました。そこで，化学構造式を\LaTeX 文書中に挿入する手順を
簡略化するパッケージを開発しようと考えました。

\clearpage

\section{パッケージを使う前に}

\subsection{依存ソフトウェアのインストール}

初めにコンピュータに\href{http://openbabel.org/}{Open Babel}と
\href{https://inkscape.org/ja/}{Inkscape}（または
\href{http://librsvg.sourceforge.net/download/}{librsvg}）をインストールし，
必要なコマンドラインプログラムである
\Lcode{obabel}と\Lcode{inkscape}（または\Lcode{rsvg-convert}）に
PATHを通します。
なお，生成画像の大きな余白を切り取るために\Lcode{pdfcrop}（または
\Lcode{ps2eps}）も利用しますが，これは\TeX\ Liveに収録されています。
正しくPATHが通っているかどうかは以下のコマンドによって確認できます：
UNIX系の場合：
\begin{verbatim}
$ which obabel
$ which inkscape        (if you choose inkscape)
$ which rsvg-convert    (if you choose librsvg)
\end{verbatim}

Windowsの場合：
\begin{verbatim}
> where obabel
> where inkscape        (if you choose inkscape)
> where rsvg-convert    (if you choose librsvg)
\end{verbatim}

何らかのディレクトリ名が返ってきた場合は，正しくインストールできているはずです。

\subsection{パッケージの読み込み}

プリアンブルに
\begin{verbatim}
\usepackage{chemobabel}
\end{verbatim}
と宣言して\Lpack{chemobabel}パッケージを読み込みます。
これは\Lpack{graphicx}パッケージに依存しますので，
もし(u)p\LaTeX を使っていてDVIドライバに\Lcode{dvipdfmx}を使用する場合は
ドライバオプション \verb+[dvipdfmx]+ をクラスオプション（グローバル）に
付けるのを忘れないようにしてください。

\clearpage

\section{基本的な使い方}
\label{basic}

\subsection{化学構造式ファイルからの変換}

\verb|ATP.cdx| はATP（アデノシン三リン酸）の構造式を
ChemDrawで描画したファイル，\verb|Brevetoxin A.mol| は
\href{http://www.chemspider.com/}{ChemSpider}から入手できるファイルです
（もとのファイル名は \verb|9041092.mol| となっています）。

\verb|\chemobabel| コマンドは本文中で以下のようにして使います。
ここでは元のファイルを\Lcode{draw}というサブディレクトリに置いたので，
ファイル名の前に \Lcode{draw/} を付けました。
もちろん \LaTeX ソースと同一ディレクトリに置くこともできます。

\begin{verbatim}
\chemobabel[width=90mm]{draw/ATP.cdx}{}
\chemobabel*[width=120mm]{draw/Brevetoxin A.mol}
\end{verbatim}

\LaTeX ファイル名がtest.texだとすると，以下のように
\verb|-shell-escape| オプションを与えてタイプセットします。
このオプションを与えることで，\LaTeX が外部プログラムである
obabelやinkscapeを呼んで実行できるようになります。
\begin{verbatim}
$ platex -shell-escape test.tex
\end{verbatim}
するとサブディレクトリ \verb|chemobabelimgdir| 以下に
「chemobabelimg[NUM].pdf」というファイルが生成します。
得られたtest.dviに対して
\begin{verbatim}
$ dvipdfmx test.dvi
\end{verbatim}
と実行すると，出力される「test.pdf」は図\ref{fig:chem1}のようになります。

\begin{figure}[h]
  \centering
  \chemobabel[width=90mm]{draw/ATP.cdx}{} \\ \vspace{1cm}
  \chemobabel*[width=120mm]{draw/Brevetoxin A.mol}
  \caption{Conversion from ATP.cdx and Brevetoxin A.mol}\label{fig:chem1}
\end{figure}

構文は以下のとおりです。
\begin{center}
\verb|\chemobabel[|\Lmeta{options}\verb|]{|\Lmeta{filename}\verb|}{|\Lmeta{obabel options}\verb|}| \\
\verb|\chemobabel*[|\Lmeta{options}\verb|]{|\Lmeta{filename}\verb|}|
\end{center}
\verb|\chemobabel| は1つのオプション引数と2つの必須引数をとります。
\begin{itemize}
 \item 角括弧内のオプション引数は \verb|\includegraphics| に渡されます。
   省略された場合は \verb|clip,scale=1| とみなされます。
 \item 必須引数の1つめには化学構造式が描かれたファイル名を与えます。
 \item 必須引数の2つめには \verb|obabel| に送るオプション
   （例は第\ref{depict}節を参照）を指定します。
\end{itemize}
必須引数の2つめを空にしたい場合は，
短縮形として \verb|\chemobabel*| を使用することもできます
（\Lver{v0.9k [2022/10/09]}以降）。

\subsection{SMILES表記法からの変換}

\LaTeX ソースファイル中に直接SMILES表記を書き込んで，
構造式に変換することもできます。
例えば，\verb|CCO| はエタノールのSMILES表記，
\verb|CC(=O)Nc1ccc(cc1)O| はアセトアミノフェンのSMILES表記です。
\verb|\smilesobabel| コマンドを以下のようにして使います：

\begin{verbatim}
\smilesobabel{CCO}{}
\smilesobabel*[width=25mm]{CC(=O)Nc1ccc(cc1)O}
\end{verbatim}

やはり \verb|-shell-escape| オプションを付けてタイプセットします。
\begin{verbatim}
$ platex -shell-escape test.tex
$ dvipdfmx test.dvi
\end{verbatim}
すると今度は「smilesobabelimg[NUM].pdf」のようなファイルが生成し，
出力される「test.pdf」は図\ref{fig:smi1}のようになります。
\begin{figure}[h]
  \centering
  \smilesobabel{CCO}{} \quad
  \smilesobabel*[width=25mm]{CC(=O)Nc1ccc(cc1)O}
  \caption{Conversion from SMILES notation}\label{fig:smi1}
\end{figure}

構文は以下のとおりです。
\begin{center}
\verb|\smilesobabel[|\Lmeta{options}\verb|]{|\Lmeta{SMILES notation}\verb|}{|\Lmeta{obabel options}\verb|}| \\
\verb|\smilesobabel*[|\Lmeta{options}\verb|]{|\Lmeta{SMILES notation}\verb|}|
\end{center}
\verb|\chemobabel| との唯一の違いは，中括弧の1つめにSMILES表記法の
テキストを与えることだけです。
\verb|\smilesobabel| の引数としては，任意のSMILES表記を用いることができます
（バックスラッシュ (\verb|\|) やパーセント記号 (\verb|%|) のような
特別な文字が入っていてもかまいません）。
なお，\verb|\chemobabel| の引数についても同様に，
特別な文字を含むことができる設計になっています。
\begin{verbatim}
\smilesobabel*[width=28mm]{Cl/C=C/Br}
\smilesobabel*[width=28mm]{Cl/C=C\Br}
\end{verbatim}
\begin{figure}[h]
  \centering
  \smilesobabel*[width=28mm]{Cl/C=C/Br} \quad
  \smilesobabel*[width=28mm]{Cl/C=C\Br}
\end{figure}


\noindent \textbf{\Lver{v0.6 [2015/06/29]}以前の旧バージョンのユーザに対する注意}

古いバージョンの\Lpack{chemobabel}では，
これらの\LaTeX で特別な意味を持つ文字に注意を払う必要がありました。
従来私が紹介してきた回避策は“カジュアルな \LaTeX ユーザ”に
カテゴリーコードを手動で変更させていたため，好ましいものではありませんでした。

この問題を解決するため，\Lver{v0.7 [2015/08/26]}以降では
パッケージ本体の側でこれらの文字をサポートするように改良を施しました。
これにより，\verb|\| や \verb|%| といった特別な文字を心配することなく，
任意のSMILES表記を直接引数に渡すことができるようになっています。
従来の回避策（\verb|\begingroup| ... \verb|\endgroup| 内で \verb|\catcode| を
変更する方法）は今ではむしろ有害になっていますので，注意してください。

\clearpage

\subsection{パッケージオプション}

\Lpack{chemobabel} パッケージは，デフォルトでは化学構造式をInkscapeによって
PDF形式の画像ファイルに変換して文書中に取り込みます。
しかし，オプションで画像形式と変換プログラムを変更することができます。

dvipsなどのPDF非対応のドライバを使用する場合は，EPS形式の画像ファイルに
変換したほうが便利でしょう。この場合
\begin{verbatim}
\usepackage[eps]{chemobabel}
\end{verbatim}
というように \verb|eps| オプションを指定すれば，内部でInkscapeから出力される
画像形式をPDFではなくEPSに変更します（\Lver{v0.9d [2016/02/28]}以降）
\footnote{\Lver{v0.9l [2022/10/22]}以降ではさらに，dvipsの場合に
自動的にEPSに切り替えるようにしました。}。

また，画像変換プログラムを Inkscape から rsvg-convert に変更することもできます
（\Lver{v0.9e [2016/03/07]}以降）。
rsvg-convertはInkscapeより高速ですから，便利かもしれません。
この場合は，以下のように \verb|librsvg| オプションを指定します。
\begin{verbatim}
\usepackage[librsvg]{chemobabel}
\end{verbatim}
デフォルトは \verb|pdf| オプションと \verb|inkscape| オプションを
指定した場合と等価です。
\begin{verbatim}
\usepackage[pdf,inkscape]{chemobabel}
\end{verbatim}

さらに，Open Babel v2.4.0以降で生成される画像はデフォルトで余白が大きいため，
\verb|pdfcrop|（PDFの場合）あるいは\verb|ps2eps|（EPSの場合）で
クロップしています（\Lver{v0.9i [2022/09/12]}以降）。
これをやめたい場合は \verb|nocrop| オプションを指定してください。
\begin{verbatim}
\usepackage[nocrop]{chemobabel}
\end{verbatim}

\clearpage

\subsection{安全なタイプセットのために}
\label{extract}

第\ref{basic}節で述べた基本的な使い方で，ほとんどの場合問題ないでしょう。
しかし，\verb|-shell-escape| オプションを付けてタイプセットするため，
好むと好まざるとにかかわらず\LaTeX が\emph{どんな外部コマンドでも}実行
できてしまいます。つまり，知らないうちにプログラムが実行されてしまう可能性を示唆
しており，これはとりわけ自分で書いた\LaTeX ソースでない場合には危険を伴います。
この問題を回避するために，\verb|\chemobabel| と \verb|\smilesobabel| コマンド
だけを別の \LaTeX ファイルに抽出するマクロを用意しています
（参考：\TeX\ Forum \cite{OKU}）。

この場合は，パッケージを読み込む際に \verb|extract| オプションを付与してください。
\begin{verbatim}
\usepackage[extract]{chemobabel}
\end{verbatim}
タイプセットの際には通常の
\begin{verbatim}
$ platex test.tex
\end{verbatim}
というコマンドを使います。
\verb|extract| オプションを付けると，\Lpack{chemobabel}パッケージは
単に元の\LaTeX ソースから \verb|\chemobabel| と \verb|\smilesobabel| コマンド
だけを抽出するはたらきをします。
その結果，作業ディレクトリに「ChemFigFile.tex」というファイルが生成します。
これは元のソースに書かれていた \verb|\chemobabel| と \verb|\smilesobabel| コマンド
をすべて含む最小の \LaTeX ソースで，これは直接pdflatexでタイプセットできます。
\begin{verbatim}
$ pdflatex -shell-escape ChemFigFile.tex
\end{verbatim}
欧文と同じ仕様に統一したため，抽出される「ChemFigFile.tex」では
\Lpack{graphicx}パッケージに \verb+[dvipdfmx]+ オプション
を与えずに読み込んでいます。したがって，
このタイプセットだけは\Lcode{pdflatex}を用いてください。
すると，第\ref{basic}節と同様にPDF形式の図のファイルが生成しますので，
あとは \verb|extract| オプションを除去して再度
\begin{verbatim}
$ platex test.tex
$ dvipdfmx test.dvi
\end{verbatim}
と実行すれば，適切に図を取り込んだ望みどおりの「test.pdf」が得られます。
この方法を用いれば，タイプセットに要する時間を大幅に短縮することも期待できます。

\paragraph{注意：\Lcode{-output-directory}オプション}

\TeX の\Lcode{-output-directory}オプションを使用すると，
抽出ファイル「ChemFigFile.tex」はその指定先に書き出されますが，
サブディレクトリ \verb|chemobabelimgdir| と生成画像は
現在のディレクトリとなります。
従って，抽出ファイルと元のファイルを正しく処理するためには
ファイルを適切に配置する必要があります。

\clearpage

\section{高度な使い方}

\subsection{描画オプション}
\label{depict}

既に述べたとおり，\verb|\chemobabel| と \verb|\smilesobabel| の両方について，
2つめの中括弧に \Lmeta{obabel options} を与えることができます。
ここで，いくつかの有用な例を挙げておきます。
\begin{itemize}
\item
\begin{verbatim}
\smilesobabel{CC(=O)Nc1ccc(cc1)O}{-xa}
\end{verbatim}
構造式中の全ての炭素原子を省略せずに描画するコマンドを実行します：
\begin{verbatim}
$ obabel -:"CC(=O)Nc1ccc(cc1)O" -O smilesobabelimg[NUM].svg -xa
\end{verbatim}
\item
\begin{verbatim}
\smilesobabel{CC(=O)Nc1ccc(cc1)O}{-xu -xC}
\end{verbatim}
元素の着色を行わず，末端の炭素原子を明示しない構造式を出力します：
\begin{verbatim}
$ obabel -:"CC(=O)Nc1ccc(cc1)O" -O smilesobabelimg[NUM].svg -xu -xC
\end{verbatim}
\item
\begin{verbatim}
\chemobabel{ATP.cdx}{-xd}
\end{verbatim}
Open Babel 2.x系列でChemDrawの図からファイル名を取り除くには，\Lcode{-xd} オプションを付けます：
\begin{verbatim}
$ obabel -:"ATP.cdx" -O chemobabelimg[NUM].svg -xd
\end{verbatim}
\end{itemize}

\begin{figure}[ht]
  \centering
  \smilesobabel[width=21mm]{CC(=O)Nc1ccc(cc1)O}{} \hspace{1cm}
  \smilesobabel[width=25mm]{CC(=O)Nc1ccc(cc1)O}{-xa} \hspace{1cm}
  \smilesobabel[width=21mm]{CC(=O)Nc1ccc(cc1)O}{-xu}
  \caption{Acetaminophen (Paracetamol): Nothing, \Lcode{-xa}, \Lcode{-xu}}
  SMILES: \verb|CC(=O)Nc1ccc(cc1)O|
\end{figure}

\begin{figure}[ht]
  \centering
  \smilesobabel[width=25mm]{C1[C@H](C)C[C@@H](O)C1}{} \hspace{1cm}
  \smilesobabel[width=25mm]{C1[C@H](C)C[C@@H](O)C1}{-xC} \hspace{1cm}
  \smilesobabel[width=25mm]{C1[C@H](C)C[C@@H](O)C1}{-xt}
  \caption{(1\textit{S},3\textit{S})-3-Methylcyclopentanol: Nothing, \Lcode{-xC}, \Lcode{-xt}}
  SMILES: \verb|C1[C@H](C)C[C@@H](O)C1|
\end{figure}

\begin{figure}[ht]
  \centering
  \smilesobabel[width=42mm]{CCc1cccc(c1)C(=O)Cl}{} \hspace{1cm}
  \smilesobabel[width=42mm]{CCc1cccc(c1)C(=O)Cl}{-xa -xu} \hspace{1cm}
  \smilesobabel[width=42mm]{CCc1cccc(c1)C(=O)Cl}{-xu -xC}
  \caption{3-Ethylbenzoyl chloride: Nothing, \Lcode{-xa -xu}, \Lcode{-xu -xC}}
  SMILES: \verb|CCc1cccc(c1)C(=O)Cl|
\end{figure}

\begin{figure}[ht]
  \centering
  \smilesobabel[width=42mm]{CCc1cccc(c1)C(=O)Cl}{--highlight "cC=O blue"} \hspace{1cm}
  \smilesobabel[width=42mm]{CCc1cccc(c1)C(=O)Cl}{-xA --genalias}
  \caption{3-Ethylbenzoyl chloride: \Lcode{--highlight "cC=O blue"}, \Lcode{-xA --genalias}}
  SMILES: \verb|CCc1cccc(c1)C(=O)Cl|
\end{figure}

\begin{figure}[ht]
  \centering
  \chemobabel[width=65mm]{draw/Firefly luciferin.mol}{} \hspace{1cm}
  \chemobabel[width=65mm]{draw/Firefly luciferin.mol}{-xa}
  \caption{Firefly luciferin (from MDL Molfiles): Nothing, \Lcode{-xa}}
\end{figure}

\begin{figure}[ht]
  \centering
  \chemobabel[width=70mm]{draw/ATP.cdx}{} \hspace{1cm}
  \chemobabel[width=70mm]{draw/ATP.cdx}{-xd}
  \caption{ATP (from ChemDraw files): Nothing, \Lcode{-xd}}
\end{figure}

\begin{figure}[ht]
より詳細にはOpen Babel documentationの
\href{http://openbabel.org/docs/dev/Command-line_tools/babel.html}{obabel and babel}または
\href{http://openbabel.org/docs/dev/FileFormats/SVG_2D_depiction.html}{SVG depiction (svg)}を
参照してください。
\end{figure}

\clearpage

\subsection{より複雑な構造式}

Open BabelはSMILES表記から複雑な構造式を生成することができます。

\begin{figure}[ht]
  \centering
  \smilesobabel*[width=55mm]{CN1CC[C@]23c4c5ccc(c4O[C@H]2[C@H](C=C[C@H]3[C@H]1C5)O)O}
  \caption{($-$)-Morphine}
  SMILES: \verb|CN1CC[C@]23c4c5ccc(c4O[C@H]2[C@H](C=C[C@H]3[C@H]1C5)O)O|
\end{figure}

\begin{figure}[ht]
  \centering
  \smilesobabel*[width=150mm]{C[C@@H]1C[C@H]2[C@@H](CC(=O)O2)O[C@H]3C[C@@H]4[C@H](C[C@@H]([C@@H]5[C@@H](O4)C/C=C\C[C@@H]6[C@@H](O5)C/C=C\[C@@H]7[C@@H](O6)CCC[C@@H]8[C@@](O7)(C[C@@H]9[C@@H](O8)C[C@@H]2[C@@H](O9)[C@H](C[C@H](O2)CC(=C)C=O)O)C)C)O[C@@]3(C1)C}
  \caption{Brevetoxin A}
  SMILES (Sorry but too long): \verb|C[C@@H]1C[C@H]2[C@@H](CC(=O)O2)O[C@H]3C[C@@H]4[C@H](C[C@@H]([C@@H]5[C@@H](O4)C/C=C\C[C@@H]6[C@@H](O5)C/C=C\[C@@H]7[C@@H](O6)CCC[C@@H]8[C@@](O7)(C[C@@H]9[C@@H](O8)C[C@@H]2[C@@H](O9)[C@H](C[C@H](O2)CC(=C)C=O)O)C)C)O[C@@]3(C1)C|
\end{figure}

しかし，時には結果が望ましくないことがあり，
また使用するバージョン番号によって出力が異なることがあります。
いくつかの例を示します：

\clearpage

\begin{figure}[p]
Cephalostatin-1 (Example from \href{http://en.wikipedia.org/wiki/Simplified_molecular-input_line-entry_system}{Wikipedia}): \\
\verb|C[C@@](C)(O1)C[C@@H](O)[C@@]1(O2)[C@@H](C)[C@@H]3CC=C4[C@]3(C2)C(=O)C[C@H]5[C@H]4CC[C@@H](C6)[C@]5(C)Cc(n7)c6nc(C[C@@]89(C))c7C[C@@H]8CC[C@@H]%10[C@@H]9C[C@@H](O)[C@@]%11(C)C%10=C[C@H](O%12)[C@]%11(O)[C@H](C)[C@]%12(O%13)[C@H](O)C[C@@]%13(C)CO|
\begin{center}
  \includegraphics[scale=0.85]{figures/Mac231-Cephalostatin-1.pdf}
  \caption{Cephalostatin-1 (figure from Open Babel for Mac 2.3.1)}
  \includegraphics[scale=0.5]{figures/Win232-Cephalostatin-1.pdf}
  \caption{Cephalostatin-1 (figure from Open Babel for Win 2.3.2)}
\end{center}
\end{figure}

\begin{figure}[p]
Sesamin: \verb|c1cc2c(cc1C3C4COC(C4CO3)c5ccc6c(c5)OCO6)OCO2|
\begin{center}
  \centering
  \includegraphics[scale=0.5]{figures/Mac231-sesamin-normal.pdf}
  \includegraphics[scale=0.5]{figures/Mac231-sesamin-gen2d.pdf}
  \caption{Sesamin (figures from Open Babel for Mac 2.3.1; Normal and \Lcode{--gen2d})}
  \includegraphics[scale=0.5]{figures/Win232-sesamin-normal.pdf}
  \includegraphics[scale=0.5]{figures/Win232-sesamin-gen2d.pdf}
  \caption{Sesamin (figures from Open Babel for Win 2.3.2; Normal and \Lcode{--gen2d})}
\end{center}
\end{figure}

\clearpage

また，Open BabelはSMILES表記から誤った構造式を生成することもあります。
ホタルルシフェリンの図で，SMILES表記から生成した構造式には
二重結合が一つ欠けています。
\begin{figure}[ht]
  \centering
  \includegraphics[scale=0.6]{figures/Exact-FireflyLuciferin.pdf}
  \caption{Firefly luciferin: exact structure from ChemSpider ID 4588411}
  \includegraphics[scale=0.6]{figures/Win232-FireflyLuciferin.pdf}
  \caption{Firefly luciferin?: output from Open Babel for Win 2.3.2}
  \smilesobabel*[width=66mm]{C1[C@@H](N/C(=c\2/nc3c(=CC(=O)C=C3)s2)/S1)C(=O)O}
  \caption{Firefly luciferin?: output from current Open Babel}
  SMILES: \verb|C1[C@@H](N/C(=c\2/nc3c(=CC(=O)C=C3)s2)/S1)C(=O)O|
\end{figure}

このような問題点は全て \verb|.mol| または \verb|.cdx| を
用意して \verb|\chemobabel| を利用することで回避できます。
とはいえ，ただの文字列から複雑な構造式を生成しうるということ自体，
おもしろいとは思いませんか？

\clearpage

\section{互換性に関する注意}

\subsection{\Lpack{chemobabel}パッケージの有利な点}

Noel O'Boyleさんによる \verb|\smiles| \cite{NOB1} あるいは
Jakob Lykke Andersenさんによる\Lpack{graphvizObabel.sty}が
提供する \verb+\obabel+ \cite{JLA} は以下のようにして使うことになっています：
\begin{verbatim}
%% ----- for \smiles command -----
\newcounter{smilescounter}
\setcounter{smilescounter}{1}
\newcommand{\smiles}[1]{
  \immediate\write18{obabel -:"#1" -O smilesimg\arabic{smilescounter}.png}
  \includegraphics{smilesimg\arabic{smilescounter}.png}
  \addtocounter{smilescounter}{1}
}
 ...
\smiles{CCO}
%% -----
\end{verbatim}
\begin{verbatim}
%% ----- for \obabel command -----
\usepackage{graphvizObabel}
 ...
\obabel[scale=0.6]{CCO}
%% -----
\end{verbatim}
しかし，\verb|\smiles| はラスター画像 (.png) を生成して
かつ \verb|\includegraphics| に渡すオプションを扱えず，\verb|\obabel| は
拡張子 (.pdf) が欠落していたりオプション引数が空の場合にエラーが発生したり
という理由で正しく処理できない可能性があります。
これらの問題は \verb|\smilesobabel| では解決されています。

\clearpage

\section{この方法の限界と代替案}

この方法はOpen Babelに依存して構造式を生成しており，
一切自分で構造式を描画することなく済ませることさえ可能です。
もちろんある程度はそれらの構造を修正したりカスタマイズしたりできますが，
この方法は出力の緻密な調節には不向きでしょう。
また，コンピュータによって生成される構造式は時に見た目が不自然で，
好ましくない場合もあるかもしれません。
したがって，この方法は構造式の出力を完全に制御したいという意図がなく，
構造式を簡便に挿入したい場合にのみ役に立つでしょう。
出力に満足できない場合は，\XyMTeX や\Lpack{chemfig}のようなパッケージを
代わりに使うことをご検討ください。

同様の目的を達成するために使えそうなパッケージとして，
以下のようなものもあります：
\begin{itemize}
  \item \href{http://www.ctan.org/pkg/mol2chemfig}{\Lpack{mol2chemfig}}:
    convert chemical structures from MDL molfile format
    to \Lpack{chemfig} source code
\end{itemize}

\clearpage

\section{技術情報}
\label{detail}

SMILES表記法に使われる特殊文字の一覧は，
\href{http://www.daylight.com/meetings/summerschool98/course/dave/smiles-intro.html}{SMILES Tutorial}%
によると以下のようです。
\begin{itemize}
\item Roman alphabets: \verb|A-Z|, \verb|a-z|
\item Numbers: \verb|1-10|
\item Brackets: \verb|[ ] ( )|
\item Others:
\begin{itemize}
\item \verb|*| (unspecified atomic number)
\item \verb|.| (disconnection)
\item \verb|+| (charge sign)
\item \verb|-| (single bond or charge sign)
\item \verb|=| (double bond)
\item \verb|#| (triple bond)
\item \verb|$| (quadruple bond)
\item \verb|:| (aromatic bond)
\item \verb|%| (used when more than 10 ring closures)
\item \verb|/|, \verb|\| (configuration around double bonds)
\item \verb|@| (tetrahedral chirality)
\item \verb|>| (reaction)
\end{itemize}
\end{itemize}

これらの文字が\LaTeX における特殊文字として解釈されることを防ぐ
必要がありますが，バックスラッシュ \verb|\| はデフォルトの
エスケープ文字であり，またパーセント記号 \verb|%| もコメント文字と
しての特殊な意味を持っています。
\Lpack{chemobabel} v0.7以上では，カテゴリーコードを一時的に変更することで
これらすべての文字を適切に扱うようになっています（\verb+\verb+と同様の手法）。

\clearpage

\section{更新履歴}

\begin{table}[h]
\centering
\begin{tabular}{ll}
2014/12/01 v0.1  & Made public as \Lpack{smilesobabel.sty} \\
2014/12/02 v0.2  & Add options which can be passed to obabel. \\
2014/12/07 v0.3  & Change name of package: \Lpack{chemobabel.sty} \\
                 & Images are stored in \Lcode{chemobabelimgdir}. \\
                 & Add \verb|\chemobabel| command. \\
2014/12/09 v0.4  & Fix a bug: (Thanks: Yusuke Terada) \\
                 & Extra spaces at the end of lines are removed. \\
2014/12/20 v0.5  & Add \verb|extract| option. \\
2015/06/29 v0.6  & Improve warning messages. \\
2015/08/26 v0.7  & Solve \verb|\catcode|-related problems in \verb|\smilesobabel|. \\
2015/08/27 v0.8  & Improve \verb|\smilesobabel| a little. \\
2015/08/28 v0.9  & Improve \verb|\smilesobabel|: (Thanks: ZR) \\
                 & Exclude $\varepsilon$-\TeX\ dependency. \\
2015/08/29 v0.9a & Improve \verb|\chemobabel| a little. \\
2016/01/06 v0.9b & Support Lua\TeX-0.85.0 and later versions. \\
2016/02/09 v0.9c & Fix a bug; forgotten in v0.9b. \\
2016/02/28 v0.9d & Add \verb|eps| and \verb|pdf| options. \\
2016/03/07 v0.9e & Add \verb|librsvg| and \verb|inkscape| options. \\
2016/10/23 v0.9f & Incorporate \verb|extract| option into main package. \\
2018/01/23 v0.9g & Minor refactor of code. \\
2022/09/11 v0.9h & Support Inkscape version 1.0 command line syntax. \\
2022/09/12 v0.9i & Crop images before inclusion, add \verb|nocrop| option. \\
2022/09/19 v0.9j & Check shell-escape status to show better error message. \\
                 & This is the first release on CTAN! \\
2022/10/09 v0.9k & Add \verb|\chemobabel*| and \verb|\smilesobabel*|. \\
                 & Warn if a user need to run another \LaTeX. \\
                 & Improve \verb|[extract]| to take over user-given package options. \\
2022/10/22 v0.9l & Default to EPS mode if PDF inclusion is unsupported by the driver. \\
                 & Default \verb|clip| for image inclusion. \\
\end{tabular}
\end{table}

\textbf{Important!}
\begin{itemize}
\item In Version 0.2, the number of parameters set in \verb|\smilesobabel| is changed!
\item From Version 0.3, the package name is changed to \Lpack{chemobabel.sty}.
\item From Version 0.7, the \verb|\catcode|-related workaround is rather harmful.
\end{itemize}

\clearpage

\begin{thebibliography}{9}
\bibitem{NOB1}
\href{http://baoilleach.blogspot.jp/2012/03/cheer-up-your-latex-with-smiles-support.html}{Cheer up your \LaTeX\ with SMILES support} -- Noel O'Blog
\bibitem{NOB2}
\href{http://baoilleach.blogspot.jp/2012/04/cheer-up-your-latex-with-smiles-support.html}{Cheer up your \LaTeX\ with SMILES support II} -- Noel O'Blog
\bibitem{JLA}
\href{http://imada.sdu.dk/~jlandersen/}{\LaTeX: Graphviz and OpenBabel} -- Jakob Lykke Andersen
\bibitem{OKU}
\href{http://oku.edu.mie-u.ac.jp/tex/mod/forum/discuss.php?d=1411}{文章内の画像のみを表示する方法} -- \TeX\ Forum \\
(How to extract only figures in \LaTeX\ source)
\bibitem{ACE1}
\href{http://acetaminophen.hatenablog.com/entry/2014/11/02/130624}{化学構造式を\TeX で(1)：自動化による簡単生成} -- Acetaminophen's diary \\
(Chemical Structural Formula in \LaTeX\ (1): An Easy Method by Auto-generation)
\bibitem{ACE2}
\href{http://acetaminophen.hatenablog.com/entry/2014/11/02/130624}{化学構造式を\TeX で(2)：自動化の注意点と解消法} -- Acetaminophen's diary \\
(Chemical Structural Formula in \LaTeX\ (2): Important Notes for Auto-generation and Solution)
\bibitem{ACE3}
\href{http://acetaminophen.hatenablog.com/entry/2014/11/05/135927}{化学構造式を\TeX で(3)：補足事項} -- Acetaminophen's diary \\
(Chemical Structural Formula in \LaTeX\ (3): Supplement)
\end{thebibliography}

\noindent
\textbf{Note}: \href{http://acetaminophen.hatenablog.com/}{Acetaminophen's diary} is my own blog! (Sorry, but only in Japanese)

I added a post about this package on the 8th day of \href{http://www.adventar.org/calendars/553}{\TeX\ \& \LaTeX\ Advent Calendar 2014} (in Japanese):
\begin{itemize}
\item \href{http://acetaminophen.hatenablog.com/entry/2014/12/08/053519}{誰でも簡単！ 化学構造式を\LaTeX に取り込むパッケージ} -- Acetaminophen's diary \\
(An easy way to insert chemical structural formulas into \LaTeX\ documents)
\end{itemize}

\end{document}
\endinput
%%
%% End of file `chemobabel-ja.tex'.
